% Source fingerprint: faad89a4c1f7b1727df9ae1a549eb15c192ff942cfe36e05f50b4bb4196f8991
% T05: raw TeX cells preserved; no numeric highlighting.
\begin{table}[htbp]
\centering\small\renewcommand{\arraystretch}{1.18}\setlength{\tabcolsep}{4pt}
\caption[线性拓扑中的网、序列与紧性]{线性拓扑中的网、序列与紧性。逐点评值的判据也适用于网。一般紧空间的网抽取不能直接改成序列抽取；Eberlein--Šmulian 的各方向按正文标明的证明范围使用。}
\label{b1:tab:norm-weak-weakstar}
\begin{tabularx}{\linewidth}{@{}>{\raggedright\arraybackslash}p{3.1cm}>{\raggedright\arraybackslash}p{4.7cm}>{\raggedright\arraybackslash}X@{}}
\toprule
收敛方式 & 测试判据 & 紧性与序列抽取 \\
\midrule
范数收敛 & \(\|x_n-x\|\to0\) & 在范数度量中，完备加全有界给出紧性；只有完备性不足以抽取子列 \\
弱收敛 & 每个 \(f\in X^*\) 满足 \(f(x_n)\to f(x)\) & 自反 Banach 空间的有界列可抽取弱子列；一般 Banach 情形用 Eberlein--Šmulian \\
\WeakStar{}收敛 & 每个 \(x\in X\) 满足 \(f_n(x)\to f(x)\) & Banach--Alaoglu 给出对偶闭球紧性；原空间可分时可度量化并抽取序列子列 \\
\bottomrule
\end{tabularx}
\end{table}
